% sections/02-related-work.tex — background + R1–R7 differentiation + positioning table.
\section{Background and Related Work}
\label{sec:related}

\subsection{Background}
Three technical threads meet in this paper. \emph{Code LLMs} provide the reasoning substrate;
their security-relevant evaluations began with secure-coding benchmarks~\cite{csev12023} and
matured into risk suites probing cyberattack helpfulness~\cite{csev22024,csev32024}.
\emph{Vulnerability detection} supplies B1's evaluation layer, where PrimeVul demonstrated that
label noise and temporal leakage had inflated a generation of results~\cite{primevul2024},
a critique systematized by Risse and B\"ohme~\cite{rissebohme2024}. \emph{Interactive
agent harnesses} descend from execution-feedback benchmarks such as InterCode-CTF
(pre-2024 background)~\cite{intercode2023} through NYU CTF Bench~\cite{nyuctfbench2024} and
Cybench~\cite{cybench2024}. Finally, the Model Context Protocol ecosystem --- measured at
scale for server-side insecurity~\cite{mcpfirstlook2025} and analyzed for tool-poisoning,
spec-level, and governance risks~\cite{mcppoisoning2026,mcpspec2026,mcpgovernance2025} ---
supplies the tool-integration substrate on which A2 autonomy operates.

\subsection{Systematizations We Position Against}
Table~\ref{tab:positioning} differentiates the seven closest prior efforts. R1's SLR mapped
300+ LLM$\times$security works but closed before the agentic turn~\cite{zhang2024when}; R2
covers detection-to-remediation with no exploitation axis~\cite{r2slr2024remediation}; R3 ---
the closest survey chronologically --- organizes pentesting only, treats ``co-evolution'' as an
architecture--benchmark relationship internal to that subfield, and excludes competition-grade
systems by corpus design; it is also outside our window (July 2026), so we use it as a
comparator, not migration evidence~\cite{r3survey2026pentest}. R4 is detection-centric within
software security~\cite{sheng2025csur}. R5 systematizes the security of agent systems ---
agents as attack surface --- whereas our subject is agents as operators of the vulnerability
loop~\cite{r5usenix26agentic}. R6's TOSEM SLR contributes 263 detection studies through
November 2025 but no offense axis and no industrial track~\cite{tosemslr2026}. R7, the AIxCC
systematization, is our most important predecessor: we cite-and-exceed it by placing CRSs on
an autonomy gradient, extending analysis to post-competition redeployment (B4), and adding the
validity critique its design-focused treatment does not attempt~\cite{r7sokaixcc2026}.

\begin{table}[t]
\centering\scriptsize
\caption{Positioning against prior systematizations. C=covered, P=partial, A=absent.}
\label{tab:positioning}
\resizebox{\columnwidth}{!}{%
\begin{tabular}{lcccccc}
\hline
Work & Window-end & Offense & Defense & Loop & Comp./industr. & Contamination\\
\hline
R1 \cite{zhang2024when}        & Aug'24 & P & C & A & A & P\\
R2 \cite{r2slr2024remediation} & Dec'24 & A & C & A & A & P\\
R3 \cite{r3survey2026pentest}  & Jul'26$^{\dagger}$ & C & P & P (intra-pentest) & A & A\\
R4 \cite{sheng2025csur}        & '25    & A & C & A & A & A\\
R5 \cite{r5usenix26agentic}    & '26    & P & P & A & P & A\\
R6 \cite{tosemslr2026}         & Nov'25 & A & C & A & P (repo) & P\\
R7 \cite{r7sokaixcc2026}       & Aug'26 & C & C & P & C & A\\
Ours                           & Jun'26 & C & C & C (graded) & C & C (measured+protocol)\\
\hline
\multicolumn{7}{l}{$^{\dagger}$Out-of-window comparator; discussed in related work only.}
\end{tabular}}
\end{table}

\subsection{What Remains Uncovered}
No prior systematization joins (i) an autonomy gradient applied to documented behavior across
academic, competition, and industrial streams; (ii) loop integration including post-competition
redeployment; and (iii) evaluation validity as a measured, protocol-addressable variable. That
conjunction is this paper's contribution.
